\documentclass[tikz,border=4pt]{standalone}
\usepackage{ctex}
\usepackage{pgfplots}
\pgfplotsset{compat=1.18}
\usetikzlibrary{calc, arrows.meta}

\begin{document}
% thinking-toolkit/09 数形结合：方程的解=图象交点
% 方程 x^2 = 2x+3 的解 = y=x^2 与 y=2x+3 的交点横坐标
\begin{tikzpicture}
\begin{axis}[
  axis lines=center,
  xlabel={$x$}, ylabel={$y$},
  every axis x label/.style={at={(axis cs:4.5,0)}, anchor=west},
  every axis y label/.style={at={(axis cs:0,12)}, anchor=south},
  xmin=-2.5, xmax=4.5, ymin=-2, ymax=12,
  xtick={-3,-2,-1,0,1,2,3,4},
  ytick={0,2,4,6,8,10},
  tick label style={font=\small},
  width=7.5cm, height=8cm,
  thick,
]
  % 抛物线 y=x^2
  \addplot[blue, thick, domain=-2.5:3.5, samples=80] {x^2};
  \node[blue, right, font=\small] at (axis cs:2.5,7) {$y=x^2$};

  % 直线 y=2x+3
  \addplot[red, thick, domain=-2.5:4, samples=20] {2*x+3};
  \node[red, right, font=\small] at (axis cs:2.8,9.5) {$y=2x+3$};

  % 两交点：x^2=2x+3 => x^2-2x-3=0 => (x-3)(x+1)=0 => x=-1,x=3
  \fill[black] (axis cs:-1,1) circle(2.5pt);
  \node[above left, font=\small] at (axis cs:-1,1) {$(-1,1)$};

  \fill[black] (axis cs:3,9) circle(2.5pt);
  \node[above left, font=\small] at (axis cs:3,9) {$(3,9)$};

  % 辅助线指向x轴
  \draw[gray, dashed] (axis cs:-1,1) -- (axis cs:-1,0)
    node[below, font=\small, black] {$x_1=-1$};
  \draw[gray, dashed] (axis cs:3,9) -- (axis cs:3,0)
    node[below, font=\small, black] {$x_2=3$};

  \node[below left, font=\small] at (axis cs:0,0) {$O$};
\end{axis}
% 注释
\node[font=\small, align=center] at (3.7,-0.7)
  {方程$x^2=2x+3$的解$\Leftrightarrow$两图象交点的横坐标};
\end{tikzpicture}
\end{document}
